\documentclass[10pt]{article}

\usepackage[utf8]{inputenc}

\usepackage[T1]{fontenc}

\usepackage{amsmath}

\usepackage{amsfonts}

\usepackage{amssymb}

\usepackage[version=4]{mhchem}

\usepackage{stmaryrd}



\begin{document}

МЕРОМОРФНЫЙ АБЕЛЕВ ДИФФЕРЕНЦИАЛ см. Абелев дифференциал.



METPИКА - обобщение понятия расстояния между точками евклидова пространства на множества, в которых можно ввести М. (метрические пространства). Для точек $x, y$ такого пространства М. $\rho(x, y)$ - это действительная неотрицательная функция, удовлетворяющая условиям:



\begin{enumerate}

  \item $\rho(x, y)=0$ лишь при $x=y$;



  \item $\rho(x, y)=\rho(y, x)$



  \item $\rho(x, y)+\rho(y, z) \geqslant \rho(x, z)$.



\end{enumerate}



Вид М. зависит как от самого пространства, так и от выбора системы координат в нем. Простейший пример М.- расстояние



\[

\left[\sum_{i=1}^{n}\left(x^{i}-y^{i}\right)^{2}\right]^{1 / 2}

\]



в декартовых координатах евклидова пространства. М. евклидова пространства в криволинейных координатах и М. риманова пространства определяются метрическим тензором. М. гильбертова пространства задается его нормой (или скалярным произведением). Понятие М. применяется и в тех случаях, когда не все условия 1)-3) выполнены: напр., если $\rho(x, y)=0$ не только при $x=y$, то $\rho$ называют п с е в д о м е т р и кой. Если М. не является положительно определенной, то ее называют индефинитной метрикой; физич. примером такой ситуации служит М. пространства Минковского в теории относительности. METРИКА ПРОСТРАНСТВА-ВРЕМЕНИ - основная геометрическая структура, которой наделяется пространственно-временное многообразие в специальной и общей теории относительности; определяется заданием поля симметричного ковариантного тензора 2-го ранга с отличным от нуля определителем - метрического тензора.



Метрич. тензор в специальной теории относительности имеет вид:



\[

\eta_{\mu \nu}=\operatorname{diag}(1,-1,-1,-1)

\]



(псевдоевклидов а метрика сигнатуры -2); пространственно-временное многообразие с такой метрикой называется пространством-временем Минковс ко го. В общей теории относительности вводится метрич. тензор $g_{\mu \nu}(x)$ более общего вида, удовлетворяющий, однако, требованию, чтобы в достаточно малой окрестности любой заданной пространственно-временной точки $x$ специальным выбором координат $g_{\mu v}(x)$ можно было свести к $\eta_{\mu \nu}$; такое пространство-время является псе вдо римано в ы п постранст в о сигнатуры -2 .



М. п.-в. задает квадрат интервала - «расстояния» между событиями, с к-рыми сопоставляются точки пространствавремени:



\[

d s^{2}=g_{\mu v} d x^{\mu} d x^{\nu}

\]



При преобразованиях пространственно-временных координат метрич. тензор, вообще говоря, изменяется (такие преобразования включают и переход к произвольно движущейся в каждой точке системе отсчета) так, чтобы величина $d s^{2}$ оставалась инвариантной. Существуют, однако,- преобразования, оставляющие метрич. тензор форминвариантным (преобразования изометрии), они выражают собой геометрич. симметрии пространства-времени, обусловленные физич. содержанием теории. Так, метрич. тензор пространства-времени Минковского в специальной теории относительности не изменяется при преобразованиях координат из группы Пуанкаре, включающих переносы начала отсчета пространственных координат и времени, повороты пространственных осей и Лоренца преобразования. Поскольку последние интерпретируются как описывающие переход от одной инерционной системы отсчета к другой, инвариантность метрики пространства-времени Минковского означает, что уравнения, записанные в лоренц-ковариантной форме, будут автоматически удовлетворять принципу относительности Эйнштейна.



В обшей теории относительности сушествование преобразований. не изменяюших М. п.-в., возможно лишь при наличии соответствуюших симметрий гравитационного піоля. Так, метрич. тензор пространства-времени Шварцшильда инвариантен относительно пространственных поворотов и временных сдвигов, что отражает центральный характер гравитационного поля и его статичность; структура метрич. тензора в моделях Фридмана, описывающих крупномасштабную структуру пространства-времени Вселенной в целом, отражает факт однородности и изотропии Вселенной в больших масштабах. Если нек-рое преобразование изометрии порождается векторным полем, то такое векторное поле называется п о л е м К и л л и нга (W. Killing, 1892) и удовлетворяет уравнению



\[

\xi_{\mu ; v}+\xi_{v ; \mu}=0

\]



где точкой с запятой обозначена ковариантная производная, согласованная с метрикой.



Следует иметь в виду, что М. П.-в. отражает не только характер гравитационного поля, но и выбор системы координат в пространстве-времени (системы отсчета). Так, переход к криволинейным координатам в пространстве-времени Минковского (к ускоренной системе отсчета) приводит к метрич. тензору общего вида, однако собственно гравитационного поля в этом случае нет. Истинное гравитационное поле связано с тензором кривизны Римана Кристоффеля, К-рый равен нулю в плоском пространстве-времени в любой системе отсчета.



М. П.-в. в случае слабого гравитационного поля непосредственно связана с ньютоновским гравитационным потенциалом $\varphi$, а именно:



\[

g_{\mu \nu} \approx \eta_{\mu \nu}+h_{\mu \nu}

\]



где $h_{\mu \nu}$ - малые добавки, характеризуюшие отклонение метрики от плоской, причем



\[

h_{00}=2 \varphi / c^{2} \text {. }

\]



Помимо задания расстояний в пространстве-времени М. п.-в. служит для определения «длины» 4-векторов $A$ :



\[

\sqrt{A^{2}}=\sqrt{A^{\mu} A^{v} g_{\mu \nu}}

\]



а также позволяет ввести операции поднятия и опускания индексов у векторов и тензоров. Определитель метрич. тензора задает инвариантный элемент объема в пространстве-времени:



\[

d \Omega=\sqrt{-g} d^{4} x,

\]



где $g$ - определитель метрич. тензора.



Д. В. Гальцов.



G-METPИКА - см. Пространство с индефинитной метрикой.



МЕТРИЧЕСКАЯ СВЯЗНОСТЬ - сМ. СвязНость.



МЕТРИЧЕСКАЯ ТЕОРИЯ ГРАВИТАЦИИ - ТеорИЯ гравитации, основанная на следующих двух принципах:



\begin{enumerate}

  \item пространственно-временное многообразие наделено метрикой;



  \item пространственно-временное многообразие эквивалентно пространству Минковского (метрика удовлетворяет принципу эквивалентности).



\end{enumerate}



Различие известных М. т. г. проявляется в законах образования метрики. Так, в общей теории относительности метрика, отличная от плоской, создается тензором энергии-импульса вещества и негравитационных полей (возможны, однако, случаи ненулевой кривизны и в пустоте), а в теории Бранса - Дикке (см. [1]) - тензором энергии-импульса вещества, негравитационных полей и постоянно присутствующего скалярного поля.



Jum.: [1] Brans G., Dicke R., «Phys. Rev.», 1961, v. 124, № 3, p. 925-35; [2] Мизнер Ч., ТорнК., Уилер Дж., Гравитация, т. 3, пер. с англ., М., 1977. А. Т. Ильичев. MЕТРИЧЕСКАЯ ТРАНЗИТИВНОСТЬ - сМ. Эргодичность.



МЕТРИЧЕСКИЙ ИНВАРИАНТ - см. Изоморфизм дИНамических систем.



МЕТРИЧЕСКИЙ 357





\end{document}